\label{sec:proximal-md-objective}

幂均值重心定义了 target，但尚未指定对固定教师最小化什么函数。本节从冻结信号积分出一个不依赖 anchor 的教师目标，证明有限步幂均值恰为其标准熵镜像下降（Mirror Descent，MD）。固定 $\kappa,\eta$ 和教师时目标不变；改变 $\eta$ 会同时改变目标与原更新步长。

\subsection{由冻结信号反求生成函数}

\begin{theorem}[幂均值信号的可积性与精确 MD 目标]
\label{thm:power-md-integrable-objective}
记 $A=A_{\kappa,\eta}$，定义规范化生成函数
\begin{equation}
 \omega_{\kappa,\eta}(x)
 =x\int_1^x\frac{A_{\kappa,\eta}(u)}{u^2}\,du,
 \qquad x>0,
 \label{eq:power-md-generator-integral}
\end{equation}
以及教师目标
\begin{equation}
 \mathcal J_{\kappa,\eta}(p;T)
 =D_{\omega_{\kappa,\eta}}(T\Vert p)
 =\sum_a p_a\omega_{\kappa,\eta}(T_a/p_a).
 \label{eq:power-md-fixed-objective}
\end{equation}
则 $\omega(1)=\omega'(1)=0$、$\omega''(x)>0$，$\mathcal J\ge0$ 且唯一零点为 $p=T$。在正正交象限内，$\mathcal J$ 的坐标梯度恰为
\begin{equation}
 \frac{\partial\mathcal J}{\partial p_a}(p;T)
 =-A_{\kappa,\eta}(T_a/p_a).
 \label{eq:power-md-exact-gradient}
\end{equation}
因此步长 $\eta$ 的 entropy MD 更新
\begin{equation}
 p^+=\argmin_{p\in\Delta}
 \left\{\langle\nabla\mathcal J(q;T),p-q\rangle
       +\frac1\eta\KL(p\Vert q)\right\}
 \label{eq:power-md-exact-step}
\end{equation}
严格等于幂均值 target $m_{\kappa,\eta}$。
\end{theorem}

\begin{proof}
对积分定义求导，得到
\begin{align}
 \omega'(x)&=\int_1^x\frac{A(u)}{u^2}\,du+\frac{A(x)}x,\\
 \omega(x)-x\omega'(x)&=-A(x),\\
 \omega''(x)&=\frac{A'(x)}x
 =\frac{x^{\kappa-2}}{1-\eta+\eta x^\kappa}>0.
 \label{eq:power-md-curvature}
\end{align}
最后一个式子对 $\kappa=0$ 按连续端点取值为 $1/x^2$。$A(1)=0$ 保证两项规范化条件。严格凸性给出 $\omega\ge0$，Jensen 等号条件和分布归一化给出散度的唯一零点。逐项使用链式法则，令 $x=T_a/p_a$，则
\begin{equation}
 \frac\partial{\partial p_a}[p_a\omega(T_a/p_a)]
 =\omega(x)-x\omega'(x)=-A(x).
\end{equation}
这一步把固定 anchor 上的信号提升为一个对任意当前分布 $p$ 定义的目标梯度。将它代入 MD 子问题，再对 $\sum p=1$ 引入乘子 $\nu$，有
\begin{equation}
 -A(T_a/q_a)+\eta^{-1}[\log(p_a/q_a)+1]+\nu=0.
\end{equation}
归一化后 $p_a^+\propto q_a\exp[\eta A(T_a/q_a)]$，正是前文的幂均值 target。该子问题严格凸，解唯一。

反过来，若规范化生成函数满足 $\omega-x\omega'=-A$，则 $(\omega/x)'=A/x^2$，连同 $\omega(1)=0$ 唯一确定式~\eqref{eq:power-md-generator-integral}。因此所构造的规范化函数是这一逐坐标梯度对应的唯一生成函数。
\end{proof}

\subsection{单调概率比信号的通用 $f$ 散度框架}

\begin{proposition}[单调信号的可积性与相对光滑性]
\label{prop:power-general-signal-framework}
设 $S:(0,\infty)\to\mathbb R$ 连续可微，$S(1)=0$、$S'(x)>0$。定义
\begin{equation}
 \omega_S(x)=x\int_1^x\frac{S(u)}{u^2}\,du,\qquad
 \mathcal J_S(p;T)=D_{\omega_S}(T\Vert p).
\end{equation}
则 $\omega_S$ 严格凸、规范化，且 $\nabla_p\mathcal J_S=-S(T/p)$。对该目标执行步长 $\tau$ 的熵 MD，精确解为 $p^+\propto q\exp[\tau S(T/q)]$。若
\begin{equation}
 L_S=\sup_{x>0}xS'(x)<\infty,
\end{equation}
则 $\mathcal J_S$ 相对于负熵 $H$ 为 $L_S$ 相对光滑；$0<\tau\le1/L_S$ 时具有本节单步下降与三点进展的相应保证。
\end{proposition}

\begin{proof}
积分微分给出 $\omega_S-x\omega_S'=-S$、$\omega_S''=S'/x>0$ 及 $\omega_S(1)=\omega_S'(1)=0$。透视函数求导后得到 $\nabla_p\mathcal J_S=-S(T/p)$，MD 一阶条件与归一化给出所列解。再令 $x_a=T_a/p_a$，有
\begin{equation}
 \nabla_p^2\mathcal J_S
 =\operatorname{diag}[x_aS'(x_a)/p_a]
 \preceq L_S\operatorname{diag}(1/p_a)=L_S\nabla^2H.
\end{equation}
沿连接两个正分布的线段积分 Hessian，得到相对光滑不等式。本节下降与三点证明只使用这个 Hessian 界、目标凸性和 MD 一阶条件，将 $1/\eta$ 替换成 $L_S$ 后逐步成立。
\end{proof}

幂均值信号属于这个框架，因为 $xA'(x)=x^\kappa/(1-\eta+\eta x^\kappa)\le1/\eta$。其他单调信号也可定义不同的凸教师目标；第~\ref{cor:power-general-signal-all-f}~推论进一步证明所列步幅下的共同教师散度进展。统一尾部界、可估计中心和良好的 BPO 投影仍需分别证明，可积性本身不保证这些额外性质。

\subsection{幂散度、reverse KL 与归一化 JS 的关系}

\begin{theorem}[精确 MD 目标的端点与曲率插值]
\label{thm:power-md-objective-endpoints}
固定正支持分布时，
\begin{align}
 \mathcal J_{0,\eta}(p;T)&=\KL(p\Vert T),\\
 \lim_{\eta\to0}\mathcal J_{\kappa,\eta}(p;T)&=D_\kappa(T\Vert p),\\
 \lim_{\eta\to1}\mathcal J_{\kappa,\eta}(p;T)&=\KL(p\Vert T).
 \label{eq:power-md-objective-limits}
\end{align}
算术形状 $\kappa=1$ 的有限 $\eta$ 目标为
\begin{equation}
 \boxed{\mathcal J_{1,\eta}(p;T)
 =\frac{(1-\eta)\KL(p\Vert M)+\eta\KL(T\Vert M)}{\eta(1-\eta)},
 \quad M=(1-\eta)p+\eta T.}
 \label{eq:power-md-js-endpoint}
\end{equation}
分子是权重为 $(1-\eta,\eta)$ 的 Jensen--Shannon（JS）散度。对所有 $x>0$，生成函数的曲率满足
\begin{equation}
 \frac1{\omega_{\kappa,\eta}''(x)}
 =\frac{1-\eta}{g_\kappa''(x)}+\frac\eta{g_0''(x)}.
 \label{eq:power-md-harmonic-curvature}
\end{equation}
所以精确 MD 目标以曲率的加权调和平均连接幂散度与 reverse KL；它一般不同于两种散度的加权和。
\end{theorem}

\begin{proof}
$\kappa=0$ 时 $A=\log x$，积分直接得到 $\omega=x-1-\log x=g_0$。在任意固定有限积分区间，$\eta\to0$ 时 $A\to(x^\kappa-1)/\kappa$，$\eta\to1$ 时 $A\to\log x$，而分母在该区间内有统一正下界。故可交换积分与极限。小 $\eta$ 时的极限生成函数与 $g_\kappa$ 有相同的二阶导数 $x^{\kappa-2}$ 及相同规范化条件；两者相等。有限词表允许逐项取散度极限。

对 $\kappa=1$，记 $a=1-\eta$、$B=a+\eta x$。分部积分给出
\begin{align}
 \omega_{1,\eta}(x)
 &=-\eta^{-1}\log B+x\int_1^x\frac{du}{u(a+\eta u)}\\
 &=\frac{x\log x}{a}-\left(\frac1\eta+\frac xa\right)\log B
 =\frac{\eta x\log x-B\log B}{\eta a}.
\end{align}
代入 $x=T/p$ 得到
\begin{equation}
 \mathcal J_{1,\eta}(p;T)
 =\frac{\eta\KL(T\Vert p)-\KL(M\Vert p)}{\eta a}.
\end{equation}
展开两项 KL，使用 $M=ap+\eta T$，其分子等于 $a\KL(p\Vert M)+\eta\KL(T\Vert M)$，得到式~\eqref{eq:power-md-js-endpoint}。最后，式~\eqref{eq:power-md-curvature} 的倒数为 $a x^{2-\kappa}+\eta x^2$；两项分别是 $a/g_\kappa''$ 与 $\eta/g_0''$。
\end{proof}

因此涉及两个 KL 方向的准确表述是：重心散度 $D_\kappa$ 在 $\kappa=0,1$ 分别为 reverse/forward KL；精确 MD 目标 $\mathcal J_{\kappa,\eta}$ 在固定正 $\eta$ 下的 $\kappa=1$ 端点是归一化加权 JS，随后令 $\eta\to0$ 才恢复 forward KL。中心 $K_{1,\eta}$ 则是 reverse KL 到混合分布，属于第三个对象。

对固定 $\kappa>0$，两个不同的插值强度一般也不能仅用目标的整体缩放互相替代。若 $\omega_{\kappa,\eta_1}=b\omega_{\kappa,\eta_2}$，则一处的二阶导数均为一，强制 $b=1$；式~\eqref{eq:power-md-curvature} 随后要求 $1-\eta_1+\eta_1x^\kappa=1-\eta_2+\eta_2x^\kappa$ 对所有 $x$ 成立，只能有 $\eta_1=\eta_2$。因此 $\eta$ 是教师目标的形状选择，同时也是原幂均值更新的 MD 步长；若要在同一个 $\mathcal J$ 上单独调节 MD 步长，应固定 $\eta$，另用后文的 $\tau$。

\begin{center}
\small
\begin{tabular}{lll}
\toprule
散度及用途 & $\kappa=0$ & $\kappa=1$、固定 $0<\eta<1$\\
\midrule
$D_\kappa(T\Vert p)$：定义重心 & $\KL(p\Vert T)$ & $\KL(T\Vert p)$\\
$\mathcal J_{\kappa,\eta}(p;T)$：精确 MD 目标 & $\KL(p\Vert T)$ & $\mathrm{JS}_\eta(p,T)/[\eta(1-\eta)]$\\
$K_{\kappa,\eta}(q,T)$：冻结教师中心 & $\KL(q\Vert T)$ & $\KL(q\Vert(1-\eta)q+\eta T)/\eta$\\
\bottomrule
\end{tabular}
\end{center}

\begin{theorem}[精确 MD 目标的锐统一上界]
\label{thm:power-md-uniform-cap}
若 $0<\kappa\le1$、$0<\eta<1$，定义
\begin{equation}
 c=-\frac{\log(1-\eta)}{\kappa\eta},\qquad
 d=\int_1^\infty\frac{A_{\kappa,\eta}(u)}{u^2}\,du.
\end{equation}
则 $0<d\le1/\eta$，并有
\begin{equation}
 0\le\mathcal J_{\kappa,\eta}(p;T)\le c+d.
 \label{eq:power-md-uniform-cap}
\end{equation}
上界不依赖词表维数，并且在至少二元的完整支持分布类上是锐上确界。算术形状满足
\begin{equation}
 c+d=\frac{-\eta\log\eta-(1-\eta)\log(1-\eta)}{\eta(1-\eta)}.
\end{equation}
原始 reverse KL 的 $\kappa=0$ 端点没有有限统一上界。
\end{theorem}

\begin{proof}
当 $u\ge1$，$A(u)>0$（除 $u=1$ 外），且 $1-\eta+\eta u^\kappa\le u^\kappa$，所以 $A(u)\le\log u/\eta$。又有 $\int_1^\infty(\log u)/u^2\,du=1$，得到 $0<d\le1/\eta$。积分定义给出 $\omega(x)/x\to d$（$x\to\infty$）。当 $x\to0$ 时 $A(x)\to-c$，写成
\begin{equation}
 \omega(x)=-x\int_x^1\frac{A(u)}{u^2}\,du\longrightarrow c.
\end{equation}
要核实最后一个极限，可对任意 $\epsilon>0$ 选取足够小的 $b$，使 $|A(u)+c|\le\epsilon$ 对 $u\le b$ 成立；区间 $[b,1]$ 的积分乘以 $x$ 趋零，区间 $[x,b]$ 的误差不超过 $\epsilon x(1/x-1/b)\le\epsilon$。

凸函数的导数不超过其无穷远渐近斜率 $d$，从零处积分可得 $\omega(x)\le c+dx$。对 $x=T/p$ 求 $p$ 期望，得到 $\mathcal J\le c+d$。取 $p=(\epsilon,1-\epsilon)$、$T=(1-\epsilon,\epsilon)$，则稀有 $p$ 坐标的贡献趋于 $d$，另一坐标的贡献趋于 $c$，证明锐性。算术形状的显式生成函数给出 $d=-\log\eta/(1-\eta)$，与 $c$ 相加得到归一化二元熵。$\kappa=0$ 时 $\omega=g_0$ 在零处发散，前述锐界退化为无穷大。
\end{proof}
\subsection{全局相对光滑、单步下降与外层收敛率}

\begin{theorem}[精确幂均值 MD 的下降与收敛]
\label{thm:power-md-relative-smoothness}
固定教师、$\kappa,\eta$，令 $H(p)=\sum_ap_a\log p_a$。在正单纯形内，$\mathcal J=\mathcal J_{\kappa,\eta}$ 相对于 $H$ 具有常数 $1/\eta$ 的相对光滑性：对任意正分布 $p,q$，
\begin{equation}
 \mathcal J(p;T)\le\mathcal J(q;T)+\langle\nabla\mathcal J(q;T),p-q\rangle
       +\eta^{-1}\KL(p\Vert q).
 \label{eq:power-md-relative-smoothness}
\end{equation}
若 $p^+\propto q\exp[\tau A(T/q)]$、$0<\tau\le\eta$，则
\begin{align}
 \mathcal J(q;T)-\mathcal J(p^+;T)
 &\ge\tau^{-1}\KL(q\Vert p^+)
  +(\tau^{-1}-\eta^{-1})\KL(p^+\Vert q),\\
 \tau\mathcal J(p^+;T)
 &\le\KL(T\Vert q)-\KL(T\Vert p^+).
 \label{eq:power-md-single-step-progress}
\end{align}
对每步刷新 anchor 的精确迭代 $q_{n+1}=p^+(q_n)$、步长 $0<\tau_n\le\eta$，有
\begin{equation}
 \mathcal J(q_N;T)\le
 \frac{\KL(T\Vert q_0)}{\sum_{n=0}^{N-1}\tau_n}.
 \label{eq:power-md-last-iterate-rate}
\end{equation}
特别地，$\tau_n=\eta$ 时每步就是幂均值 target，最后迭代误差不超过 $\KL(T\Vert q_0)/(\eta N)$。
\end{theorem}

\begin{proof}
记 $x_a=T_a/p_a$。根据式~\eqref{eq:power-md-exact-gradient}，目标在正正交象限的 Hessian 为
\begin{equation}
 \nabla_p^2\mathcal J(p;T)
 =\operatorname{diag}\left(\frac1{p_a}
       \frac{x_a^\kappa}{1-\eta+\eta x_a^\kappa}\right)
 \preceq\eta^{-1}\operatorname{diag}(1/p_a)
 =\eta^{-1}\nabla^2H(p).
 \label{eq:power-md-hessian-bound}
\end{equation}
对连接 $q,p$ 的线段积分二阶 Taylor 余项，得到相对光滑不等式。目标本身严格凸，其 Hessian 的每个对角元为正。

设 $g=\nabla\mathcal J(q;T)$。MD 的一阶条件给出 $g_a+\tau^{-1}\log(p_a^+/q_a)$ 为常数。与 $p^+-q$ 内积时常数消失，故
\begin{equation}
 \langle g,p^+-q\rangle
 =-\tau^{-1}[\KL(p^+\Vert q)+\KL(q\Vert p^+)].
\end{equation}
代入相对光滑式，得到第一个下降界。再用凸性 $\mathcal J(q;T)-\mathcal J(T;T)\le\langle g,q-T\rangle$，以及 $\mathcal J(T;T)=0$，得到
\begin{align}
 \mathcal J(p^+;T)
 &\le\langle g,p^+-T\rangle+\eta^{-1}\KL(p^+\Vert q)\\
 &=\tau^{-1}[\KL(T\Vert q)-\KL(T\Vert p^+)]
   -(\tau^{-1}-\eta^{-1})\KL(p^+\Vert q)\\
 &\le\tau^{-1}[\KL(T\Vert q)-\KL(T\Vert p^+)].
\end{align}
第二行使用熵 Bregman 的三点恒等式，可直接展开各 KL 核验。对 $n=0,\ldots,N-1$ 求和，非负的末项 KL 给出
\begin{equation}
 \sum_{n=0}^{N-1}\tau_n\mathcal J(q_{n+1};T)\le\KL(T\Vert q_0).
\end{equation}
由第一个下降界，$\mathcal J(q_N;T)\le\mathcal J(q_{n+1};T)$，故左侧至少为 $(\sum\tau_n)\mathcal J(q_N;T)$，得到最后迭代界。
\end{proof}

这个定理无需概率统一正下界，只要求给定有限分布严格为正；初始 $\KL(T\Vert q_0)$ 可以很大。若 $\sum\tau_n\to\infty$，则 $\mathcal J(q_N;T)\to0$，严格凸散度及有限单纯形紧性保证 $q_N\to T$。边界也不能形成额外的零点：由积分定义，
\begin{equation}
 \lim_{x\to\infty}\frac{\omega(x)}x
 =\int_1^\infty\frac{A(u)}{u^2}\,du=:d_\infty>0.
\end{equation}
该积分有限，因为 $0\le A(u)\le\log u/\eta$ 对 $u\ge1$ 成立。若某个坐标 $p_a\to0$，则 $p_a\omega(T_a/p_a)\to T_a d_\infty>0$。任何内点极限的零损失又强制各比率都等于一。因此所有收敛子列的极限只能是 $T$，整个序列收敛到 $T$。

对固定 anchor 的 BPO 内层，理想路径 $p_\xi\propto q\exp[\xi\eta A(T/q)]$、$0\le\xi\le1$ 对应从该 anchor 做一步步长为 $\tau=\xi\eta$ 的上述 MD；因而具有同一单步界。这里的外层收敛率要求每步重新计算信号并刷新 anchor。实际共享参数 BPO/AdamW 只有近似投影时，需要另外控制拟合误差；任意 BPO loss 下降不能自动套用该定理。若外层改变 $\eta$ 或 $\kappa$，目标 $\mathcal J$ 也改变，不能把各步的不等式直接望远镜求和。
